\section*{Köttbullar}
\ihead{}\chead{Personen:2}\ohead{}
\ifoot{Vorbereitungszeit:35 min}\cfoot{}\ofoot{Kochzeit:25 min}
{\Large Zutaten}
\begin{multicols}{2}
\begin{itemize}
\item \SI{500}{g} Rinderhackfleisch oder Elchhackfleisch
\item \num{1} Zwiebel
\item \num{1} Ei, mittelgroß
\item \SI{2}{EL} Paniermehl
\item \SI{1}{TL} Sojasauce
\item \SI{1}{TL} Salz
\item \SI{1/2}{TL} Pfeffer
\item \SI{1}{TL} Gewürzmischung Piffi Allkrydda, ähnlich Steak-Grill-Gewürz
\item \SI{50}{g} Butter zum Braten
\item \SI{200}{g} Schlagsahne
\item \SI{100}{ml} Wasser
\item \SI{1}{TL} Mehl
\item \num{1} Prise Zucker
\item Salz
\item Pfeffer
\item Gewürzmischung Piffi Allkrydda
\item \num{1} Glas Wildpreiselbeeren
\item \SI{500}{g} mehligkochende Kartoffeln
\item Butter, optional
\end{itemize}
% \columnbreak
\end{multicols}
\noindent {\Large Zubereitung}\\
\\
\textit{Vorbereitung:} Das Hackfleisch in eine mittelgroße Schüssel geben.
Die Zwiebel schälen, in feine Würfel schneiden und zum Hackfleisch geben.
Ei, Paniermehl und Gewürze hinzufügen und alles miteinander vermengen.
Die Hände befeuchten und aus der Hackfleischmasse kleine Bällchen formen.
Je kleiner die Bällchen, desto besser.
Den Ofen auf \SI{100}{\celsius} Ober-/Unterhitze vorheizen.
Die Butter in einer Pfanne erhitzen und die Bällchen von allen Seiten scharf anbraten.
Wenn sie durchgebraten sind, in eine Schüssel geben und im Ofen warmhalten.
\textit{Soße:} Das Fett in der Pfanne mit dem Wasser ablöschen.
Unter Rühren die Sahne hineingießen und aufkochen lassen.
Das Mehl mit etwas Wasser verrühren und die Sahnesoße damit eindicken.
Mit Salz, Pfeffer und Piffi Allkrydda abschmecken.
\textit{Beilage:} Die Kartoffeln schälen, vierteln und in kochendem Salzwasser ca. \SI{20}{min} garen.
Anschließend abgießen und je nach Geschmack ein kleines Stück Butter zu den Kartoffeln geben.
Den Topf wieder auf den Herd stellen und die Herdplatte ausschalten.
Die Köttbullar aus dem Ofen nehmen, aber nicht in die Soße legen.
Alles separat mit den Wildpreiselbeeren servieren.
